\documentclass[aspectratio=169]{beamer}

\usetheme{Madrid}
\usecolortheme{default}
\setbeamertemplate{navigation symbols}{}

\usepackage{amsmath}
\usepackage{booktabs}
\usepackage{array}

\title[VokVision Interim]{VokVision: Interim Presentation}
\subtitle{Monocular Multi-view 3D Reconstruction Pipeline}
\author{Abhigyan Raj \and Mantavya Kaushik \and Shubhi Jain}
\institute{Engineering Track}
\date{April 2026}

\begin{document}

% Slide 1
\begin{frame}
  \titlepage
\end{frame}

% Slide 2
\begin{frame}{1) Problem Statement, Scope, and Users}
\textbf{Problem}\
Build an end-to-end system where users capture multi-view RGB photos and receive a downloadable 3D model.

\vspace{0.5em}
\textbf{Input}\
Monocular image set (JPEG/PNG) captured from multiple viewpoints.

\vspace{0.5em}
\textbf{Output}\
Reconstructed artifact exported as \texttt{point\_cloud.ply} (and Gaussian model artifacts).

\vspace{0.5em}
\textbf{Target Users}\
Students, creators, and rapid-prototyping teams needing capture-to-3D without specialized sensors.

\vspace{0.5em}
\textbf{Current Scope}\
\begin{itemize}
  \item Mobile capture + backend asynchronous reconstruction.
  \item Monocular-only; quality depends on viewpoint coverage.
  \item Object isolation works with automatic fallback if SAM2 is unavailable.
\end{itemize}
\end{frame}

% Slide 3
\begin{frame}{2) Related Work and Baselines (I)}
\textbf{Classical Geometry Baseline: COLMAP-style SfM}
\begin{itemize}
  \item Strength: robust camera pose + sparse geometry on many datasets.
  \item Method gap: feature matching can fail on low-texture / repetitive regions.
  \item Method gap: hand-crafted matching + geometric verification can be brittle under motion blur.
\end{itemize}

\vspace{0.5em}
\textbf{Neural View Synthesis Baseline: NeRF}
\begin{itemize}
  \item Strength: high-fidelity novel-view synthesis.
  \item Method gap: slow training/inference for interactive product workflows.
  \item Method gap: optimization can be unstable when capture viewpoints are sparse/non-uniform.
\end{itemize}
\end{frame}

% Slide 4
\begin{frame}{2) Related Work and Baselines (II)}
\textbf{Modern 3D Priors: DUSt3R / MASt3R}
\begin{itemize}
  \item Key idea: infer geometry/matching directly in 3D rather than only 2D descriptors.
  \item Strength: better scene-level geometric grounding for challenging correspondence.
  \item Method gap: quality still sensitive to image coverage and camera path diversity.
\end{itemize}

\vspace{0.5em}
\textbf{Rendering Baseline: 3D Gaussian Splatting}
\begin{itemize}
  \item Key idea: explicit Gaussian primitives for faster rendering than NeRF.
  \item Strength: near real-time rendering after training.
  \item Method gap: densification/opacity heuristics can be scene-sensitive and error-prone.
\end{itemize}

\vspace{0.5em}
\textbf{Chosen baselines for this project}
\begin{itemize}
  \item COLMAP-style SfM
  \item NeRF-style neural rendering
  \item 3D Gaussian Splatting target pipeline
\end{itemize}
\end{frame}

% Slide 5
\begin{frame}{3) Dataset and Evaluation Metrics}
\textbf{Dataset snapshot (current interim experiments)}
\begin{itemize}
  \item Latest validated run: 22 JPEG images (project: \texttt{test\_project\_001}).
  \item Workspace assets: \texttt{gs\_dataset/images} and \texttt{gs\_dataset/sparse}.
  \item Practical issue: small internal set, limited diversity across devices/scenes.
\end{itemize}

\vspace{0.4em}
\textbf{Evaluation Metrics}
\begin{itemize}
  \item Reconstruction quality: PSNR, SSIM, LPIPS.
  \item Geometry quality: Chamfer distance / completeness (when GT available).
  \item System metrics: end-to-end runtime, GPU VRAM, CPU RAM, storage footprint.
  \item Product metrics: successful job rate, average time-to-result, API error rate.
\end{itemize}

\vspace{0.4em}
\textbf{Resource relevance}
GPU memory and storage throughput are the dominant constraints.
\end{frame}

% Slide 6
\begin{frame}{4) System Overview and Current Implementation}
\textbf{Current end-to-end pipeline}
\begin{enumerate}
  \item Flutter app captures and uploads image set.
  \item FastAPI receives files and starts background job.
  \item Stage 1: MASt3R mapping (poses + geometry).
  \item Stage 2: object point-cloud isolation (SAM2 or fallback) + dataset prep.
  \item Stage 3: OpenSplat training and model export.
\end{enumerate}

\vspace{0.5em}
\textbf{Completed components}
\begin{itemize}
  \item Upload/status/download API flow.
  \item Pipeline shell with stage orchestration.
  \item Integration of MASt3R, isolation stage, and OpenSplat codebases.
\end{itemize}
\end{frame}

% Slide 7
\begin{frame}{4) Baseline Results and Status Analysis}
\textbf{Module-wise interim status}
\begin{table}[h]
\centering
\begin{tabular}{p{2.8cm}p{4.8cm}}
\toprule
Module & Interim status \\
\midrule
Local orchestration & Functional (22 images, 2500 iterations) \\
MASt3R Stage & Completed (2,181,208 pts written) \\
Isolation Stage & Completed (final 8,294 clean pts) \\
OpenSplat Stage & Completed (2500 steps) \\
Output artifact & Generated: \texttt{model.ply} (934.2 KB) \\
\bottomrule
\end{tabular}
\end{table}

\textbf{Observation}
Current bottlenecks are environment/dependency robustness and benchmark depth, not core pipeline completion.
\end{frame}

% Slide 8
\begin{frame}{4) Error Analysis, Failure Cases, and Insights}
\textbf{Observed failure modes}
\begin{itemize}
  \item VLM branch skipped because \texttt{GEMINI\_API\_KEY} was missing.
  \item SAM2 module unavailable, forcing GrabCut+HSV fallback segmentation.
  \item MASt3R used slow RoPE2D PyTorch fallback (missing CUDA-compiled extension).
  \item HF Hub requests unauthenticated (no \texttt{HF\_TOKEN}), affecting reliability under load.
\end{itemize}

\vspace{0.5em}
\textbf{Likely causes (method + system)}
\begin{itemize}
  \item Methodology: segmentation quality is sensitive when fallback replaces learned masks.
  \item System: optional runtime dependencies and tokens are not validated pre-run.
  \item System: coarse status reporting still slows root-cause diagnosis.
\end{itemize}

\vspace{0.5em}
\textbf{Engineering implication}
Need strict preflight checks, richer telemetry, and reproducible environment locking.
\end{frame}

% Slide 9
\begin{frame}{5) Individual Contributions So Far}
\textbf{Abhigyan Raj}
\begin{itemize}
  \item Core reconstruction integration work and pipeline-level debugging.
\end{itemize}

\textbf{Mantavya Kaushik}
\begin{itemize}
  \item Backend/API orchestration and stage execution path setup.
\end{itemize}

\textbf{Shubhi Jain}
\begin{itemize}
  \item Workflow validation, reporting, and integration/testing coordination.
\end{itemize}

\vspace{0.5em}
\textbf{Team-level completed milestone}
Interim end-to-end architecture assembled, with clear blocker localization.
\end{frame}

% Slide 10
\begin{frame}{5-6) Next Steps, Task Assignment, and Demo Plan}
\textbf{Next steps (near-term)}
\begin{itemize}
  \item Wire this successful local config into API-triggered production path.
  \item Add explicit job states: queued/running/failed/succeeded.
  \item Add metric logging for runtime, memory, and quality.
  \item Add preflight checks for keys/modules (GEMINI\_API\_KEY, HF\_TOKEN, SAM2).
  \item Expand dataset diversity for stronger EDA and stress tests.
\end{itemize}

\vspace{0.5em}
\textbf{Future ownership}
\begin{itemize}
  \item Abhigyan Raj: reconstruction-core fixes + quality analysis.
  \item Mantavya Kaushik: backend reliability + status/error API enrichment.
  \item Shubhi Jain: experimental protocol, result synthesis, and presentation/demo flow.
\end{itemize}

\vspace{0.5em}
\textbf{Demo plan}
Run a 22-image job, show stage-wise logs, and download/render the generated \texttt{model.ply}.
\end{frame}

% Backup slides
\appendix

\begin{frame}{Backup: Compute Requirements}
\begin{itemize}
  \item GPU: CUDA-capable, preferably at least 12GB VRAM.
  \item CPU RAM: 16--32GB recommended.
  \item Storage: SSD required for I/O-heavy training workflow.
  \item Operational need: queue-based scheduling for multi-user load.
\end{itemize}
\end{frame}

\begin{frame}{Backup: Known Technical Debt}
\begin{itemize}
  \item Optional dependencies/tokens are not enforced before run start.
  \item Coarse status model (no explicit failed state).
  \item Deprecated package/runtime warnings (\texttt{google.generativeai}, Python 3.10 horizon).
  \item Missing CUDA-compiled RoPE2D extension causes slower MASt3R execution.
\end{itemize}
\end{frame}

\end{document}
